\chapter{Introducción}

\section{Objetivos}

El presente trabajo práctico tiene como objetivo diseñar y analizar un filtro pasa alto pasivo de segundo
orden con topología RLC, para la frecuencia de corte asignada al grupo, considerando dos valores del factor
de calidad $Q$. En particular, se busca:

\begin{itemize}
  \item Obtener la función de transferencia del filtro en función de los valores de $R$, $L$ y $C$, y
        dimensionar los componentes para cada valor de $Q$.
  \item Calcular las raíces del polinomio característico y ubicarlas en el plano complejo.
  \item Calcular el factor de amortiguamiento e identificar el tipo de amortiguamiento de cada filtro.
  \item Obtener y graficar la respuesta temporal ante un escalón, midiendo sobre ella los parámetros
        característicos (sobrepico, tiempo de caída y tiempo de establecimiento).
  \item Comparar ambos diseños y describir las ventajas de cada tipo de amortiguamiento.
\end{itemize}

\section{Descripción del trabajo}

Los parámetros de diseño corresponden al Grupo 5 de la Tabla~1 de la consigna y se resumen en la
Tabla~\ref{tab:parametros}. El valor de la resistencia $R$ es una elección de diseño y se justifica en el
Capítulo~\ref{chap:diseno}.

\begin{table}[!ht]
  \centering
  \begin{tabular}{|l|c|}
    \hline
    Parámetro & Valor \\
    \hline
    Frecuencia de corte $f_c$ & \qty{400}{\hertz} \\
    Factor de calidad $Q_1$   & \num{0.25} \\
    Factor de calidad $Q_2$   & \num{0.75} \\
    Resistencia $R$ (elegida) & \qty{100}{\ohm} \\
    \hline
  \end{tabular}
  \caption{Parámetros de diseño asignados al grupo.}
  \label{tab:parametros}
\end{table}

El desarrollo se organiza de la siguiente manera: en el Capítulo~\ref{chap:diseno} se plantea la topología,
se deduce la función de transferencia y se obtienen las ecuaciones de diseño comunes a ambos casos. En los
Capítulos~\ref{chap:q025} y~\ref{chap:q075} se analiza cada filtro por separado (función de transferencia
numérica, raíces, amortiguamiento y respuesta temporal). En el Capítulo~\ref{chap:comparacion} se comparan
ambos diseños y se discuten las ventajas de cada tipo de amortiguamiento. Finalmente, el
Capítulo~\ref{chap:conclusiones} presenta las conclusiones. Todos los resultados analíticos fueron
verificados numéricamente con el script que se adjunta en el Anexo~\ref{chap:anexo}.
